\hypertarget{appendix-d67a19507dc5}{}
\section*{Machine word}
\section{Machine word}\label{machine-word}

\textbf{Curriculum depth:} Developed · \textbf{First seen:} C001 · \textbf{Latest development:} C002

\subsection{Definition and mechanism}\label{definition-and-mechanism}

A context-dependent natural processing width, not a portable fixed-size unit.

A processor's natural integer processing width is often called its word width, but the term is architecture dependent. A 64-bit CPU still supports smaller integer operations and wider vector operations. State exact widths in formats instead of relying on the word label.

\subsection{Prerequisites}\label{prerequisites}

\begin{itemize}
\tightlist
\item
  \hyperref[appendix-5bee2df85484]{Byte and octet} --- An octet is exactly eight bits.
\end{itemize}

\subsection{Used by}\label{used-by}

\begin{itemize}
\tightlist
\item
  \hyperref[appendix-8f38927cfc54]{Bitwise NOT} --- uses machine word as part of its representation or reasoning.
\item
  \hyperref[appendix-73425b25118e]{Shift operations} --- uses machine word as part of its representation or reasoning.
\item
  \hyperref[appendix-c99977d88b33]{Unsigned modular arithmetic} --- uses machine word as part of its representation or reasoning.
\item
  \hyperref[appendix-e12d1b03fefb]{Atomic operations and write contention} --- uses machine word as part of its representation or reasoning.
\item
  \hyperref[appendix-942ad41afb30]{Compiler IR and machine instructions} --- uses machine word as part of its representation or reasoning.
\item
  \href{}{SIMD and GPU data layouts (02-Foundations-II, appendix-3c7b1a56de01)} --- uses machine word as part of its representation or reasoning.
\item
  \hyperref[appendix-cc04044b5782]{Fixed-width integer types} --- uses machine word as part of its representation or reasoning.
\end{itemize}

\subsection{Recall and next depth}\label{recall-and-next-depth}

Explain machine word without the label, work the example above, then state the assumption whose removal would change the result.

\textbf{Pending:} registers, instruction width, SIMD lanes, and ABI data models.

\subsection{Source and navigation}\label{source-and-navigation}

\href{Continuous\%20Notes/Computing\%20Foundations\%20-\%20From\%20Bits\%20to\%20C++\%20and\%20Systems\%20APIs\#From\%20a\%20signal\%20to\%20an\%20interpretation}{C001 teaching chapter} · \href{Continuous\%20Notes/Computing\%20Foundations\%20-\%20From\%20Bits\%20to\%20C++\%20and\%20Systems\%20APIs\#Arithmetic\%20becomes\%20a\%20policy}{C002 development} · \hyperref[appendix-fd090d982a15]{Part definitions and mechanisms} · \href{}{Dependency Map (02-Foundations-II, appendix-4ad47bdd269a)}
